\documentclass[11pt]{article}
\usepackage[a4paper,margin=1in]{geometry}
\usepackage{amsmath,amssymb,amsthm,hyperref}
\newtheorem{theorem}{Theorem}
\newtheorem{proposition}{Proposition}
\newtheorem{corollary}{Corollary}
\title{GST All-Graph Causal Holonomy, Reciprocal-Fold Hodge Birth, and Unbounded Purity Gaps}
\date{22 October 2026}
\begin{document}
\maketitle
\begin{abstract}
We derive a genuinely geometric all-graph extension of the GST dimensionless-world and higher-causal-fold framework. Every original smooth projective complex variety X enters through B=X^m with its actual cyclic algebraic factor automorphisms. Twisted graph incidence, not an abstract Hodge-class operator, computes the emergent weight-defective cohomology. Multiple loops produce source sectors even with zero common holonomy invariants. Two reciprocal loops supply an exact full-target cohomological retraction, and arbitrarily deep causal tensor folds produce every negative purity gap. For finite cyclic holonomy, a smooth projective companion has an exact pure-Hodge polynomial. We distinguish this structure from any claimed construction of a previously unknown algebraic cycle on X.
\end{abstract}

\section{Geometric GST input: worlds before dimensions}
Fix a connected smooth projective complex variety X, dim_C X=d, and an arbitrary integer m>=2. Let B=X^m, b=md, with algebraic cyclic factor permutation rho; on rational cohomology set T=rho^*. This is the actual realization of a causal world automorphism: no metric, dimension, cohomological degree or Hodge type is primitive in the underlying GST GeneralSpace, but a chosen algebraic realization produces them.

Let Gamma be a finite connected graph with r vertices and e edges, loops and parallel edges permitted, and g=e-r+1 independent cycles. Assign to every oriented edge a an algebraic automorphism f_a of B, with inverse orientation having inverse f_a, and assume all f_a generate a finite group. For each vertex take B x P^1, choose different marked points for incident half-edges, and glue the two B-sections for each edge through f_a. Denote the resulting projective pinched scheme Y. It has complex dimension b+1; normalization consists of r disjoint B x P^1; double locus consists of e disjoint B. Projectivity follows by equipping an ample tensor-orbit bundle of the finite holonomy group with its group linearization, tensoring it with O_{P^1}(n), descending along marked sections, and using finite-surjective descent of ampleness. This is an actual topology-changing geometric realization.

\section{Twisted incidence law in every Hodge degree}
Write V_q=H^q(B,Q), a pure polarizable rational Hodge structure of weight q. Set C^0_q=V_q^{\oplus r}, C^1_q=V_q^{\oplus e}. The twisted graph coboundary is
\[
 (\delta_q x)_a = T_a x_{t(a)}-x_{s(a)}, \qquad T_a=f_a^*.
 \tag{1}
\]
Its kernel K_q and cokernel C_q are genuine rational Hodge structures because each T_a comes from an algebraic automorphism.
\begin{theorem}[All-graph causal mixed-Hodge birth]
For every cohomological degree k,
\[
\boxed{\operatorname{Gr}_{k-1}^{W}H^k(Y,Q)=C_{k-1},\qquad
\operatorname{Gr}_{k}^{W}H^k(Y,Q)=K_k\oplus V_{k-2}(-1)^{\oplus r}.}\tag{2}
\]
Here V_j=0 for j<0, and all other weight-graded pieces vanish. These are associated graded isomorphisms; a canonical splitting of the entire mixed Hodge structure is NOT asserted.
\end{theorem}
\begin{proof}
For the normalization nu:tilde Y -> Y and disjoint double locus N, the rational constant-sheaf normalization sequence
0 -> Q_Y -> nu_* Q_tildeY -> Q_N -> 0
maps branch values to their twisted differences. Under Kunneth, H^q(tildeY)=V_q^{\oplus r} \oplus V_{q-2}(-1)^{\oplus r} and H^q(N)=V_q^{\oplus e}; the connecting map is delta_q on the first summand and zero on the second, because O(1) on P^1 restricts trivially to a marked point. Its long exact sequence gives
0 -> C_{k-1} -> H^k(Y) -> K_k \oplus V_{k-2}(-1)^{\oplus r} -> 0.
The first term is pure of weight k-1 and the last pure of weight k. Strictness of the mixed-Hodge weight filtration proves (2).
\end{proof}

\begin{theorem}[All-graph causal source amplification]
The twisted spanning-tree gauge identifies K_q with the simultaneous fixed subspace V_q^{Hol(Gamma)} of the cycle holonomies. Rank-nullity yields
\[
\boxed{\dim C_q=(g-1)\dim V_q+\dim V_q^{Hol(Gamma)}.}\tag{3}
\]
Thus for any g>=2, at least (g-1)dim V_q independent weight-defective causal channels exist even if every common holonomy-fixed direction vanishes.
\end{theorem}
\begin{proof}
The tree coordinates identify all kernel vertex vectors with a single vector fixed by each fundamental cycle holonomy. If beta_q=dim V_q and R_q=rank delta_q, then dim K_q=r beta_q-R_q and dim C_q=e beta_q-R_q. Subtraction gives dim C_q=(e-r) beta_q+dim K_q=(g-1)beta_q+dim K_q.
\end{proof}

\section{Exact all-dimension Hodge-generating equation}
Define the degree/weight/Hodge polynomial P_B(t,s,u,v)=sum h^{p,q}(Gr^W_w H^k(B,Q))t^k s^w u^p v^q, and A=t^2 s^2 uv. Let P_K denote the Hodge polynomial of all holonomy-invariant cohomology. The exact sequence defining C_q implies its pure-Hodge Grothendieck class [C_q]=(g-1)[V_q]+[K_q]. Therefore
\[
\boxed{P_Y=(1+t)P_K+[rA+(g-1)t]P_B.}\tag{4}
\]
This recovers a tree of r rational components as (1+rA)P_B and the prior one-loop cyclic GST formula when g=r=1. For cyclic factor permutation of B=X^m, the exact graded Koszul-signed character average is
\[
\boxed{P_K=\frac1m\sum_{\ell\mid m}\varphi(\ell)
\left[P_X((-1)^{\ell-1}t^\ell,s^\ell,u^\ell,v^\ell)\right]^{m/\ell}.}\tag{5}
\]
The signs are necessary when H^odd(X) is nonzero.

\section{Reciprocal causal loops give a full-Hodge retraction}
On a one-vertex graph take two loop holonomies T and T^{-1}, both induced by genuine algebraic automorphisms. Their first two incidence coordinates are
\[
 \Delta(v)=((T-I)v,(T^{-1}-I)v).
\]
Define R(a,b)=a+Tb. The exact identity
\[
\boxed{R\Delta=(T-I)+T(T^{-1}-I)=0}\tag{6}
\]
shows that R descends to a surjective Hodge map barR:C_q -> V_q. The insertion S(v)=[(v,0)] is a section,
\[
\boxed{\bar R S=I_{V_q}.}\tag{7}
\]
This survives any additional graph loops because R ignores their coordinates. It does not require a nonzero holonomy-invariant vector.

For every original X, integer p and rational Hodge class alpha in H^{2p}(X,Q) intersect H^{p,p}(X), choose B=X^m and T=rho^*. The point-inclusion i:X -> B, x -> (x,x0,...,x0), satisfies i^* pr_1^*=id. Hence
\[
\boxed{\Hdg^p_Q(X)\xrightarrow{pr_1^*}V_{2p}
 \xrightarrow{S}C_{2p}
 \cong Gr^W_{2p}H^{2p+1}(Y,Q)}
 \tag{8}
\]
is injective, with exact return i^* barR S pr_1^*(alpha)=alpha. The target alpha need NOT be algebraic, and no Hodge-to-Chow inverse has been postulated. This improves the earlier cyclic symmetrization by avoiding its averaging step.

\section{All causal depths and exact purity signatures}
Let D be a genuine irreducible nodal rational projective curve, whose H^1(D,Q)=Q(0) is in degree one and weight zero. For every h>=1 set
Z_h=Y x D^{h-1}, an actual projective singular geometry of complex dimension md+h. Then
\[
\boxed{P_{Z_h}=P_Y(1+t+A)^{h-1}.}\tag{9}
\]
Mixed-Hodge Kunneth applied to the independent weight-zero graph holes gives for EVERY original-X rational Hodge class alpha
\[
\boxed{\Hdg^p_Q(X)\hookrightarrow
Gr^W_{2p}H^{2p+h}(Z_h,Q),\quad
(p',q')=(p,p),\ \chi=p'-q'=0,\ \delta=(p'+q')-(2p+h)=-h.}
\tag{10}
\]
The causal depth h is unbounded. For even h the born rational diagonal sector lies in even cohomological degree but has the wrong purity weight for the classical smooth-projective Hodge target. Z_h legitimately has a mixed Hodge structure, not a counterexample to the Hodge conjecture on a smooth projective variety.

\section{Smooth projective pure-Hodge companion for cyclic holonomy}
Assume g>=1 and the g holonomies generate a finite cyclic group G. Choose a smooth projective genus-g curve C, and map its fundamental group generators a_i to the chosen holonomies, b_i to identity; this defines a surjection to G, hence a connected finite étale G-cover C' -> C. The free diagonal quotient S=(B x C')/G is genuinely smooth projective, of dimension b+1. It is a companion geometric realization, not an automatically established flat smoothing of the singular Y.

For the nontrivial character line local systems of C' -> C, Riemann--Roch gives holomorphic one-form multiplicity g-1; the trivial character contributes g. Consequently the entire pure Hodge polynomial is
\[
\boxed{P_S=(1+A)P_K+ts(u+v)[(g-1)P_B+P_K].}\tag{11}
\]
Both source and smooth companion depend on the same B cohomology and finite cyclic holonomy, but different realized topologies; every smooth-fiber Hodge sector has weight equal to degree. A smoothness claim for Y, or a generic Hodge-surjectivity claim on X, is not inferred.

\section{Adversarial geometric example: ZERO fixed classes, nonzero births}
Let E be any elliptic curve and B=E x E. On V=H^1(B,Q)=Q^4 let S exchange factors and [-1] be simultaneous elliptic negation. The two distinct involutions f_1=S and f_2=[-1]S act as T_1=S and T_2=-S on V. Their common fixed subspace is 0. Yet the two-loop incidence matrix Delta:Q^4 -> Q^8 has rank 4, hence coker dimension 4. With P_\pm=(I \pm S)/2, define
\[
 R(a,b)=P_+a-P_-b,\quad
 \Sigma(v)=(P_+v,-P_-v).
\]
Then R Delta=0 and R Sigma=I_4 by exact rational matrix identities. This uses different algebraic automorphisms, not duplicated identity loops. The generated group is C_2 x C_2, so formula (11), proved for cyclic G, is not applied to its smooth companion without extending the representation argument.

\section{Canonical cycle-class firewall: source labels versus output topology}
Write A_p(B)=im(cl_B^{canonical}:CH^p(B)_Q -> H^{2p}(B,Q)). Algebraic holonomy maps preserve A_p(B); the reciprocal return R and injection S preserve the corresponding algebraic edge labels. Therefore
\[
\boxed{\gamma\in A_p(B)
\iff
S(\gamma)\in im(A_p(B)^2\to coker\Delta_{2p}).}
\tag{12}
\]
The reverse direction applies R to algebraic edge labels and uses R Delta=0. Applying i^* to gamma=pr_1^*alpha yields the equivalence with algebraicity on the untouched original X. Importantly, an algebraic node-label is NOT automatically a codimension-(p+1) Chow cycle class in the ordinary cohomology of a singular ambient Y: cycle class for singular schemes requires the correct Chow/Borel--Moore or operational semantics. The topology births Hodge-labelled cohomology, not a new algebraic preimage for arbitrary alpha.

\section{Executed exact computation and verification boundaries}
Wolfram rational-matrix tests of cyclic shifts in dimensions m=2,...,7 and g=2,...,5 passed all 24/24 instances of (3). With elliptic X, m=2,...,4 and g=1,...,4, all 12/12 singular/smooth Hodge-polynomial pairs had nonnegative integer coefficients. At m=2,g=2, the cyclic singular fold has weight-one H^2 polynomial 3u+3v and the smooth companion has pure degree-two polynomial 3u^2+9uv+3v^2. A separate noncyclic two-involution calculation gives ker Delta=0, dim coker=4 and R Delta=0, R Sigma=I_4, all exact. These computations are finite checks; the equations above are derived uniformly under the geometric assumptions.

Mathbox provenance: mathematically self-reviewed LaTeX derivation; scheme-level normalization/MHS, finite covers and Riemann--Roch are mathematical inputs, NOT claims of current Lean-kernel verification. No local Lean installation, no fabricated Hodge source and no public web search. The original all-X rational Hodge conjecture remains a separate exact target: producing an arbitrary original-X native rational Chow cycle with the genuine canonical Betti class is not a consequence of these causal sector equations.
\end{document}